\begin{abstract}
This paper presents an empirical assessment of privacy and security risk posture in 15 Android consumer apps using ScytaleDroid, a database-backed analysis platform. Build-backed evidence was collected within a fixed 14-day cutoff window, preserving APK hashes, package versions, static and dynamic run identifiers, PCAP availability, and quality-assurance provenance. Static analysis examines permissions, component exposure, detector findings, and MASVS-aligned control areas across harvested base and split APK artifacts. Runtime analysis complements that evidence with strict-idle, quiescent foreground, and interactive captures from a non-root physical device. The results show measurable static exposure alongside runtime patterns that are not visible from packaged artifacts alone, including a near-zero association between static finding volume and interactive retained-host count. The study further shows that a recent cutoff window is a more auditable reporting gate than same-day installed-build completion under app-update churn.
\end{abstract}

\begin{IEEEkeywords}
Android security, mobile application privacy, static analysis, runtime analysis, dynamic analysis, APK analysis, MASVS, privacy risk, security risk, evidence provenance
\end{IEEEkeywords}
